\section{Discussion and Future Work}
\label{sec:conclusion}

The headline of this paper is intentionally narrow. We re-implemented
\steermoe~close to the upstream paper, applied it to \mixtral~on a
question-only fear-concept test, swept the natural knobs, and observed
small, mostly-defocusing shifts rather than concept-specific transfer.
The risk-difference signal that drives the steering plan is diffuse on
\mixtral, which is consistent with the question-only behavior.

\paragraph{What this is, and is not, evidence for.}
We are not claiming \steermoe~is broken. The original demonstration uses
different MoE checkpoints, different prompt families, and a different
test contract. We are claiming that on the specific transfer setup we
ran --- \mixtral, fear concepts, question-only test --- the upstream
recipe is not enough on its own. That is a useful diagnostic both for
practitioners considering routing-bias steering on Mixtral-class models
and for follow-up work that wants to compare readout rules at fixed
intervention site.

\paragraph{Same-data, same-model method comparison.}
The natural follow-up is a same-Mixtral, same-data comparison of three
columns:

\begin{enumerate}
  \item \mixtral~baseline,
  \item \mixtral~+ attention-guided activation steering on hidden states,
  \item \mixtral~+~\steermoe~(this paper's setup).
\end{enumerate}

This is the experiment our codebase is staged for. The
attention-collection scaffolding under
\texttt{collect\_attention\_to\_prefix.py} and
\texttt{src/.../attention\_collection.py} already produces the
attention-to-prefix tensors that an attention-guided readout requires;
what is missing is the hidden-state readout-and-edit backend that turns
those tensors into a steering plan. Adding that backend without
changing the dataset, the prompt contracts, or the report schema is a
deliberate design choice of the current repo.

\paragraph{Concept-family generalization.}
A second axis we did not explore is whether the diffuse-signal finding
holds on more lexical concept families. If a political-figure or
profession concept family produces a sharply concentrated
risk-difference table on the same Mixtral checkpoint, the diffuse-signal
explanation would localize to the fear-concept setting rather than to
\mixtral. This would shift the implication of Claim~3 from
``\steermoe-on-\mixtral~needs a better readout'' toward
``\steermoe-on-\mixtral~needs a more lexical target concept''. We
consider both worth running before any method-comparison conclusion.

\paragraph{What we will and will not try to claim later.}
The codebase keeps the model-specific and method-specific pieces
deliberately separated --- \texttt{mixtral\_backend.py} is the active
router-tracing layer, \texttt{olmoe\_backend.py} is the previous
implementation we used to validate the readout shape, and the
attention-collection module is shared. That separation is what lets
this paper make a narrow Claim~2/3 about \steermoe-on-\mixtral~without
overcommitting to a method comparison the data cannot yet support.
The next paper, if the planned method comparison comes through, will
have to be more ambitious; this one does not.
